\section{Modalità di misurazione}
Per effettuare le misurazioni sul tempo di risoluzione del sistema, abbiamo utilizzato le funzionalità offerte nativamente dai
rispettivi linguaggi:
\begin{itemize}
    \item \textbf{Matlab}: abbiamo utilizzato i comandi \texttt{tic} e \texttt{toc} per rispettivamente far partire il timer
    e fermarlo.
    \item \textbf{C++}: per effettuare la misurazione, il programma C++ calcola la differenza di tempo da poco dopo aver letto la matrice in memoria
    a subito dopo la risoluzione del sistema. Per ottenere i due istanti di tempo necessari al calcolo, abbiamo utilizzato la funzione \texttt{std::chrono::high\_resolution\_clock::now()}
    della libreria \texttt{chrono} del linguaggio.
\end{itemize}

\lstinputlisting[
    language=C++,
    style = Cpp-editor,
    firstline = 105,
    lastline = 133,
    firstnumber = 105,
    caption = Porzione del codice C++ in cui si calcola il tempo di esecuzione
]{../cpp/utils/utils.hpp}
\clearpage
\lstinputlisting[
    style=Matlab-editor,
    firstline=66,
    lastline=72,
    firstnumber=66,
    caption = Porzione del codice Matlab in cui si calcola il tempo di esecuzione
]{../matlab/main.m}
\section{Risultati}




